\documentclass[10pt,a4paper]{amsart}
\usepackage[margin=19mm]{geometry}
\usepackage{amsmath,amssymb,mathtools}
\usepackage[scr=rsfs,cal=euler]{mathalfa}
\usepackage{xfrac}
\usepackage[unicode,hidelinks,bookmarks=false]{hyperref}
\newcommand{\ie}{i.e.\ }
\newcommand{\cf}{cf.\ }
\newcommand{\resp}{respectively\ }
\newcommand{\Subsection}{\subsection}
\providecommand{\nobreakdash}{\nobreak\hbox{-}}
\setlength{\emergencystretch}{2em}
\multlinegap0pt
\usepackage{bm}
\usepackage{bbm}
\usepackage{dsfont}
\usepackage{xspace}
\usepackage{cancel}
\usepackage{mathtools}
\usepackage[shortlabels]{enumitem}
\usepackage{amsthm}
\makeatletter
\@ifundefined{theorem}{\newtheorem{theorem}{Theorem}}{}
\@ifundefined{lemma}{\newtheorem{lemma}[theorem]{Lemma}}{}
\@ifundefined{observation}{\newtheorem{observation}[theorem]{Observation}}{}
\makeatother
\newcommand{\T}{\mathcal{T}}
\newcommand{\eps}{\varepsilon}
\newcommand{\ff}{\mathcal{F}}
\newcommand{\g}{\mathcal{G}}
\newcommand{\s}{\mathcal{S}}
\renewcommand{\ge}{\geqslant}
\renewcommand{\le}{\leqslant}
\title[A complete $t$-intersection theorem for families of spanning]{A complete $t$-intersection theorem for families of spanning trees}
\author{Elizaveta Iarovikova, Andrey Kupavskii}
\date{Source: arXiv:2507.17913v1 (CC BY 4.0)}
\begin{document}
\maketitle

\noindent\emph{Source and licence.} A complete $t$-intersection theorem for families of spanning trees. Version arXiv:2507.17913v1 (CC BY 4.0), distributed under the Creative Commons Attribution 4.0 International licence, as recorded in the arXiv licence metadata for this exact version and shown on the abstract page https://arxiv.org/abs/2507.17913v1. Full rights notice: licenses/arxiv-2507.17913-CC-BY-4.0.txt.\par\smallskip

\noindent\textit{Editorial scope.} Counted unit: the Theorem whose source label is \texttt{1SA} in the article (source0.tex lines 318--326, source label \texttt{1SA}), delivered together with the article's own complete original proof of it (source0.tex lines 328--364, one \texttt{proof} environment of 37 lines). No mathematics is written, repaired or re-derived: statement and proof are the article's own text line for line. Required context retained uncounted: obs\_spread\_2 (lines 146--148); spread\_lemma (lines 152--154); cnt\_relative (lines 214--222). Every definition, display and label of those lines is retained unchanged. Omitted from this package and not redistributed: 1--145, 149--151, 155--213, 223--317, 327--327, 365--666. Nothing inside a retained excerpt is omitted or altered: every retained body is delivered exactly as the article's lines are written, so no omission receipt of any kind accompanies this package. Citations: the retained excerpts carry no citation command at all, so no bibliography entry of the article is transcribed and no reference is redistributed. Reference adaptation: none was needed and none was made; every reference inside the delivered unit resolves inside it, so no unresolved reference or citation is produced. Printed slips kept literal: no printed word, symbol, hypothesis or number of the counted statement or of its proof was altered. Layout and class: the article's own class is \texttt{amsart} with packages including cmap, fontenc, inputenc, babel, color, amsmath, amssymb, amsthm, arcs, enumitem, epsfig, filecontents, graphicx, hyperref; the wrapper is the lane's pinned \texttt{amsart} editorial wrapper at 10pt on A4, which restates the theorem-like environments the retained text needs and adds the article's own macro definitions that the retained text needs (\texttt{\textbackslash T}, \texttt{\textbackslash eps}, \texttt{\textbackslash ff}, \texttt{\textbackslash g}, \texttt{\textbackslash ge}, \texttt{\textbackslash le}, \texttt{\textbackslash s}), with extra wrapper packages (bm, bbm, dsfont, xspace, cancel, mathtools, enumitem). The delivered numbering is the unit's own. Assets: no retained span contains a figure environment, a graphics-inclusion command, a PDF-page inclusion, a table or a TikZ picture; no image file is copied into this package, the package declares no asset dependency and no image file is redistributed with it. Licence: version arXiv:2507.17913v1 of this article is distributed under the Creative Commons Attribution 4.0 International licence, as recorded in the arXiv licence metadata for this exact version and shown on the abstract page https://arxiv.org/abs/2507.17913v1. A full rights notice is delivered at licenses/arxiv-2507.17913-CC-BY-4.0.txt. Characters: every retained excerpt is pure ASCII, so the delivered engine, XeLaTeX with its default Latin Modern text font, reports no missing character.
\begin{observation}\label{obs_spread_2}
  Fix $\alpha>1$ and consider a family $\ff$. If $X$ is an inclusion-maximal set with the property that $\ff(X) \ge \alpha^{-|X|}|\ff|$, then $\ff(X)$ is $\alpha$-spread.
\end{observation}

\begin{theorem}\label{spread_lemma}
  If for some $n,k,r\ge 1$ a family $\ff\subset {[n]\choose \le k}$ is $r$-spread and $W$ is an $(\beta\delta)$-random subset of $[n]$, then $$\Pr[\exists F\in \ff\ :\ F\subset W]\ge 1-\Big(\frac 2{\log_2(r\delta)} \Big)^\beta k.$$
\end{theorem}

\begin{lemma}\label{cnt_relative}
    Let $n,t, r$ be nonnegative integers such that $t + r\le n - 2$. Then 
    \begin{enumerate}[label=(\roman*)]
        \item $c_{n, t+r}/c_{n, t} \le 2^r$;
        \item if $t \ge 2n/3$ then $c_{n, t+r}/c_{n, t} \le (e/3)^r$.
        \item if $n/2 \le t < 2n/3$ then $c_{n, t+r}/c_{n, t} \le (9/8)^r$.
        \item $c_{n, t+r}/c_{n, t} \ge 1/n^r$
    \end{enumerate}
\end{lemma}

\begin{theorem}\label{1SA}
  There is such $\eps_0 > 0$ that for any $\eps < \eps_0, n > n_0(\eps)$ the following holds. Let $t\ge1$ be an integer such that $n^{1-\eps/3} \le t \le n- 2$. Consider a family $\ff\subset \T_n$ that is $t$-intersecting.
  Put $t'= t-n^{1-\epsilon/3}$, $q' = \min\{(1 + \eps)t, n-1\}$. Then there exists a  $t'$-intersecting family $\mathcal S\subset \T_n^{\le q'}$ of spanning forests of size at most $q'$ and a family $\ff'$ such that the following holds.
  \begin{itemize}
    \item[(i)] $\ff\setminus \ff'\subset \T_n[\s]$;
    \item[(ii)] for any $B\in \s$ there is a family $\ff_B\subset \ff$ such that $\ff_B(B)$ is $n^{\epsilon/2}$-spread;
    \item[(iii)] $|\ff'|\le n^{-\frac 13\eps t}\cdot c_{n, t}n^{n-t-2}$.
  \end{itemize}
\end{theorem}

\begin{proof}[Proof of Theorem~\ref{1SA}] We construct $\mathcal S$ as follows. Put $r:=%\frac \epsilon 4
n^{\eps/2}$ and consider the following procedure for $i=1,\ldots $ with $\ff^1:=\ff$.
\begin{enumerate}
    \item Find a maximal $S_i$ that  $|\ff^i(S_i)|\ge  r^{-|S_i|}|\ff^i|$.
    \item If $|S_i|> (1+\eps)t$ or $\ff^i = \emptyset$ then stop. Otherwise, put $\ff^{i+1}:=\ff^i\setminus \ff^i[S_i]$.
\end{enumerate}
Note that Observation~\ref{obs_spread_2} and maximality of $S_i$ imply that $\ff^i(S_i)$ is $r$-spread.
Let $N$ be the step of the procedure for $\ff$ at which we stop. The family $\s$ is defined as follows: $\s:=\{S_1,\ldots, S_{N-1}\}$. Clearly, $|S_i|\le q'$ for each $i\in [N-1]$. The family $\ff_{B}$ promised in (ii) is defined to be $\ff^i[S_i]$ for $B=S_i$. Next, note that if $\ff^N$ is non-empty, then $(1+\eps)t < |S| \le n-1$. Put $a = |S| - t > \eps t$. Then
\begin{align*}
|\ff^N|\le&r^{|S_N|}  |\ff^{N}(S_N)|\le  r^{|S_N|}c_{n, |S_n|}n^{n-2-|S_n|} \le \\ 
\overset{\ref{cnt_relative}, (i)}{\le}&
n^{\frac 12\eps|S_N|}\cdot 2^a\cdot c_{n, t}\cdot n^{n - 2 - t - a} = \\
=& c_{n, t}n^{n-2-t}n^{\eps t/2 + \eps a/2 - a+a\log_n2} \le \\
\le& n^{-\frac13\eps t}\cdot c_{n, t} \cdot n^{n-2-t}. 
\end{align*}
The last inequality is because $$a\left(\frac \eps2 - 1 + \log_n2\right) + \frac{\eps t}2 \leqslant -\frac{\eps t}3 \Longleftrightarrow$$ $$\Longleftrightarrow a\ge \frac{5 \eps t}{6(1 - \eps/2 - \log_n2)}.$$ which holds for $\eps < \eps_0$ and $n > n_0(\eps)$ if we fix $\eps_0$ small enough and $n_0(\eps)$ large enough, because $$a \ge \eps t \ge  \frac{5 \eps t}{6(1 - \eps/2 - \log_n2)}.$$ We are only left to verify that $\s$ is $t'$-intersecting if we chose $\eps_0$ and $n_0(\eps)$ appropriately. We prove it below.

\begin{lemma}\label{lemtint_1SA} The family $\mathcal S$ is $t'$-intersecting. \end{lemma}
\begin{proof}

  Take two (not necessarily distinct) $A_1,A_2\in \mathcal S$  and assume that $|A_1\cap A_2|<t'$. Recall that the families $\ff_{A_i}(A_i)$, $i=1,2$ are both $r$-spread with $r = n^{\eps/2}$. 
  For $i\in [2]$ put $j:=3-i$ and put
  $$\g_i:=\ff_{A_i}(A_i)\setminus\Big\{F\in \ff_{A_i}(A_i): |F\cap A_j\setminus A_i|\ge \frac{n^{1-\eps/3}}2\Big\}.$$ The size of the latter family is at most
 \begin{align*}{|A_j\setminus A_i|\choose n^{1-\epsilon/3}/2}\max_{X: |X| = n^{1-\eps/3}/2, X\cap A_i = \emptyset}|\ff_{A_i}(A_i\cup X)|\le&\\  {n\choose n^{1-\eps/3}/2}n^{-\frac\eps 4 n^{1-\eps/3}}|\ff_{A_i}(A_i)|\le&\\
  (2e n^{\eps/3})^{\frac 12n^{1-\eps/3}}n^{-\frac\eps 4 n^{1-\eps/3}}|\ff_{A_i}(A_i)|\le&\\
  n^{\frac{\eps}5n^{1 - \eps/3} - \frac{\eps}4n^{1 - \eps/3}}|\ff_{A_i}(A_i)| \le&\frac 12 |\ff_{A_i}(A_i)|.\end{align*}
  Two last inequalities require $n$ to be large enough, the exact minimal value of $n$ depends on $\eps$ only. We can guarantee that by the choice of $n_0(\eps)$. This implies that $|\g_i|\ge \frac 12 |\ff_{A_i}(A_i)|$.
 Because of this and the trivial inclusion $\g_i(Y)\subset \ff_{A_i}(A_i\cup Y)$, valid for any $Y$,  we conclude that $\g_i$ is $ \frac r2$-spread for both $i\in [2]$. Since $n$ can be sufficiently large depending only on $\eps$, we have $\frac r2 > 2^{11}\log_2(2n)$.  



  What follows is an application of Theorem~\ref{spread_lemma}. Let us put $\beta= \log_2(2n)$ and $\delta = (2\log_2(2n))^{-1}$. Note that $\beta\delta = \frac 12$ and $\frac r2\delta > 2^{4}$ by our choice of $r$.  Theorem~\ref{spread_lemma} implies that a $\frac{1}2$-random subset $W_i$ of ${{[n]\choose 2}\choose 2}\setminus A_i$ contains a set from $\g_i$ with probability strictly bigger than
  $$1-\Big(\frac 2{\log_2 2^{4}}\Big)^{\log_2 (2n)} n = 1-2^{-\log_2 (2n)} n = \frac 12.$$

  Consider a random partition of ${[n]\choose 2}\setminus (A_1\cup A_2)$ into $2$ parts $U_1',U_2'$, where $\Pr[x\in U_i']=1/2$ independently for each $i\in [2]$ and $x\in {[n]\choose 2}\setminus (A_1\cup A_2)$. Next, for $i\in[2]$ and $j=3-i$, put $U_i:= U_i'\cup (A_j\setminus A_i)$. Then  $U_i$ is a random subset of the same ground set as $W_i$ above, moreover, the probability of containing each element is at least that for $W_i$. It implies that there is $F_i \in \g_i$ such that $F_i\subset U_i$ with probability strictly bigger than $\frac 12$. Using the union bound, we conclude that, with positive probability, it holds that there are such $F_i$, $F_i\subset U_i,$ for both  $i \in[2]$. Fix such choices of $U_i$ and $F_i$, $i \in [2]$. Then, on the one hand, each $F_i\cup A_i$ belongs to $\ff$ and, on the other hand, $|(F_1\cup A_1)\cap (F_2\cup A_2)| =|A_1\cap A_2|+|F_1\cap A_2|+|F_2\cap A_1|< |A_1\cap A_2|+n^{1-\epsilon/2}<t$. The first inequality follows from the definition of $\g_i$, and the second inequality due to $|S_1\cap S_2|<t'$ and the definition of $t'$. This is a contradiction with $\ff$ being $t$-intersecting.
  \end{proof}
This concludes the proof of Theorem~\ref{1SA}.\end{proof}

\end{document}
